% =============================================================================
%  APPENDIX C - Statistical Methods
% =============================================================================
\chapter{Statistical Methods}
\label{chap:appC}

\chapterquote{All statistical inference is a conversation between the data
and the assumptions.}{after George Box}

\begin{chapterpreview}
This appendix gives formal definitions for the three statistical
procedures used in \cref{chap:morph-results}: the Shapiro--Wilk test
for normality, the Wilcoxon signed-rank test for paired changes
when the normality assumption fails, and the Tukey
$1.5\times$\gls{IQR} criterion for outlier identification on the
volume-change distribution. The appendix also documents the
significance thresholds, the multiple-comparison correction, and the
power-analysis rationale for the chosen sample size ($n=45$ tracer
particles).
\end{chapterpreview}

\section{Notation}

In this appendix, $\{x_i\}_{i=1}^{n}$ denotes a sample of size $n$ drawn
from a population with unknown distribution. For paired data
(pre-test / post-test of the same particle), $\{d_i\}_{i=1}^{n}$
denotes the per-particle change, $d_i = x_i^{\mathrm{after}} -
x_i^{\mathrm{before}}$. The median of a sample is denoted
$\widetilde{x}$ and the order statistics by
$x_{(1)} \le x_{(2)} \le \dots \le x_{(n)}$.

\section{Shapiro--Wilk test for normality}
\label{sec:appC-shapiro}

The Shapiro--Wilk statistic $W$ tests the null hypothesis $H_0$ that a
sample is drawn from a normal distribution against the alternative
$H_1$ that it is not. The statistic is
\begin{equation}
W \;=\; \frac{\left(\sum_{i=1}^{n} a_i\, x_{(i)}\right)^{2}}
              {\sum_{i=1}^{n}\left(x_i - \bar{x}\right)^{2}},
\label{eq:appC-shapiro}
\end{equation}
where $\bar{x}$ is the sample mean and the coefficients $\{a_i\}$ are
tabulated constants derived from the expected values and the
covariance matrix of the order statistics of a standard normal sample.
The statistic lies in the closed interval $[0,1]$; values close to $1$
favour $H_0$, and small values reject it.

\paragraph{Application in this thesis.} The test is applied in
\cref{sec:morph-volume} to the paired volume-change distribution
$\{\%\Delta V_i\}$ across the $n=45$ tracer particles. A rejection of
$H_0$ ($p < 0.05$) motivates the use of a non-parametric paired test
(Wilcoxon signed-rank) in place of the parametric paired $t$-test.

\section{Wilcoxon signed-rank test for paired data}
\label{sec:appC-wilcoxon}

The Wilcoxon signed-rank statistic $T^+$ tests the null hypothesis
$H_0$ that the median of the paired-difference distribution is zero
against the alternative $H_1$ that it is non-zero, without requiring
the paired differences to be normally distributed. The procedure is:
\begin{enumerate}[leftmargin=*]
    \item Compute the per-pair differences
          $d_i = x_i^{\mathrm{after}} - x_i^{\mathrm{before}}$.
    \item Discard any pairs with $d_i = 0$ and reduce $n$ accordingly.
    \item Rank the absolute values $|d_i|$ from $1$ (smallest) to $n$
          (largest), assigning the mean rank to ties.
    \item Compute
          \begin{equation}
              T^+ \;=\; \sum_{i:\,d_i>0} R_i,
              \qquad
              T^- \;=\; \sum_{i:\,d_i<0} R_i,
              \label{eq:appC-wilcoxon}
          \end{equation}
          where $R_i$ is the rank of $|d_i|$.
    \item The smaller of $T^+$ and $T^-$ is referred to the
          Wilcoxon-signed-rank reference distribution to obtain a
          two-sided $p$-value; for $n>20$ the normal-approximation
          $z$-statistic
          \begin{equation}
              z \;=\; \frac{T^+ - n(n+1)/4}{\sqrt{n(n+1)(2n+1)/24}}
              \label{eq:appC-wilcoxon-z}
          \end{equation}
          is used.
\end{enumerate}

\paragraph{Application in this thesis.} The test is applied to the
paired sphericity, convexity, angularity, roundness, and volume changes
across all $n=45$ tracer particles. The sphericity result reported in
\cref{sec:morph-shape} is $z = -5.84$, $p < 10^{-6}$, supporting the
universal \gls{ModeA} (diffuse-abrasion) claim with overwhelming
statistical evidence.

\section{\texorpdfstring{Tukey $1.5\times$\gls{IQR} outlier criterion}
                         {Tukey 1.5xIQR outlier criterion}}
\label{sec:appC-tukey}

For a sample with inter-quartile range $\IQR = Q_3 - Q_1$, the Tukey
outlier fence places a point $x_i$ in the outlier set if either
\begin{equation}
x_i \;<\; Q_1 - 1.5 \cdot \IQR
\qquad\text{or}\qquad
x_i \;>\; Q_3 + 1.5 \cdot \IQR,
\label{eq:appC-tukey}
\end{equation}
where $Q_1$ and $Q_3$ are the first and third quartiles of the sample
respectively. The factor $1.5$ is the classical Tukey choice; under a
normal distribution, this fence excludes approximately
\SI{0.7}{\percent} of the population on each tail. The criterion is
robust because $Q_1$, $Q_3$, and $\IQR$ are themselves robust
location-and-scale estimators.

\paragraph{Application in this thesis.} The criterion is applied to
the paired volume-change distribution $\{\%\Delta V_i\}$ in
\cref{sec:morph-volume}. The $Q_1, Q_3$ values are computed from all
$n=45$ paired records; particles falling below
$Q_1 - 1.5\,\IQR$ are labelled \gls{ModeB} (catastrophic-fracture
candidates). The six identified outliers are exactly the six particles
of \cref{fig:morph-volume-scatter}.

\section{Significance thresholds and multiple-comparison correction}

The significance threshold throughout the thesis is the conventional
$\alpha = 0.05$ unless stated otherwise. Where a family of $k$
comparisons is made simultaneously (the sphericity, convexity,
angularity, and roundness Wilcoxon tests of \cref{sec:morph-shape},
$k=4$), the Bonferroni-corrected threshold $\alpha/k = 0.0125$ is
applied to control the family-wise Type-I error rate. All four
descriptors meet this stricter threshold; the principal sphericity
result ($p<10^{-6}$) is well below it.

\section{Sample-size rationale}

The sample size of $n=45$ tracer particles (five per configuration
across nine configurations) was determined \emph{a priori} from a
two-sided paired-test power analysis. With significance level
$\alpha=0.05$ and target power $1-\beta = 0.80$, a paired test on the
volume-change distribution can detect an effect size of
$|d|/\sigma_d \approx 0.42$ at $n=45$. The smallest practically
meaningful effect size in the morphological literature is
$|d|/\sigma_d \approx 0.5$ (corresponding to a sphericity change of
$\Delta \Psi \approx 0.003$ against a typical within-population
standard deviation of $\sigma_\Psi \approx 0.006$). The chosen
sample size therefore comfortably exceeds the power requirement for
the principal effect of interest.

\begin{thesisnote}[title=Reproducibility]
All statistical procedures in this appendix are implemented in the
\texttt{scipy.stats} module of the Python pipeline of
\cref{chap:appB}: \texttt{scipy.stats.shapiro},
\texttt{scipy.stats.wilcoxon} (with \texttt{zero\_method="wilcox"}
and \texttt{correction=True}), and a hand-coded Tukey fence based on
\texttt{numpy.percentile}. The pipeline produces deterministic
output from the paired \gls{STL} meshes and is therefore fully
reproducible by any reviewer with access to the raw data.
\end{thesisnote}

\begin{chaptertakeaways}
\begin{itemize}[leftmargin=*,itemsep=2pt,topsep=2pt]
    \item \textbf{Shapiro--Wilk} verifies that the
          parametric-normality assumption fails for the paired
          volume-change data; the non-parametric Wilcoxon test is
          therefore the correct inference tool.
    \item \textbf{Wilcoxon signed-rank} confirms that \gls{ModeA}
          (diffuse abrasion) is universal: the median paired
          sphericity change is significantly greater than zero,
          $p<10^{-6}$.
    \item \textbf{Tukey $1.5\times$\gls{IQR}} identifies six
          \gls{ModeB} (catastrophic-fracture) outliers in the
          volume-change distribution; all six are drawn from rigid-base
          configurations only.
    \item All three procedures are implemented in the open Python
          pipeline of \cref{chap:appB}.
\end{itemize}
\end{chaptertakeaways}
